% introduction.tex

\section*{Introduction}

Machine learning can screen molecular designs before experimental testing, but
its usefulness is limited by the labels available for training. In materials
informatics, this bottleneck has motivated simulation-to-real (Sim2Real)
learning, in which models are trained with large computational databases from
first-principles calculations, quantum chemistry, or molecular dynamics, then
adapted to smaller experimental data sets
\citep{Yamada2019Predicting,Aoki2023Multitask,Minami2025Scaling}. In some
polymer and inorganic-material settings, transfer error decreases as the
computational database grows, providing a practical test of whether additional
simulation is worth generating \citep{Minami2025Scaling}. That test is needed
because simulation data are not automatically useful. Model assumptions, finite
sampling, and domain gaps can separate a computed label from the experimental
property it is meant to support.

Protein engineering has the same label bottleneck, but the relation between
computational and experimental labels is often less direct. Experimental measurements are sparse, whereas
simulations and structure-based calculations can label many designed or mutated
sequences. Those computed labels, however, rarely measure the target phenotype
itself. For thermal stability, the experimental target may be the melting
temperature (Tm), the assay-dependent temperature at which folded and unfolded
populations are balanced. The available computed quantity may instead be a
mutation-level free-energy change ($\Delta\Delta G$), a structure-based stability
score, or a molecular-dynamics trajectory summary. Such labels are physically
related to stability, but they are not interchangeable with Tm. The central
mismatch motivates a direct test of which computed quantities transfer to
experimental Tm prediction.

We study this question for nanobodies, single-domain VHH antibodies that are
small, modular, and useful as engineered binders
\citep{Muyldermans2013Nanobodies,Alexander2024Discovery}. Their practical use in
expression, storage, and formulation depends on thermal stability
\citep{Kunz2018-qt,Mezzasalma2007-ci,Kumar2000-te}, but measuring or simulating
that stability directly is costly
\citep{Bekker2019Thermal,Higashida2021-lt,Mohseni2019-lx,Tomimoto2023-ww}.
Protein language models (PLMs) such as ESM2 provide a natural starting point
because they learn sequence representations from large protein corpora
\citep{Rives2021Biological,Lin2023Evolutionary,Hayes2025-pd}, and several recent
methods use such representations for Tm prediction
\citep{Blaabjerg2023Rapid,Li2023DeepTM,Alvarez2024TEMPRO,Savojardo2024-tv,Ramon2025-oy,Mao2026-yf}.
High-throughput stability assays are also changing the field. Tsuboyama et al.
measured folding stability for roughly 776{,}000 natural and designed protein
domains by cDNA-display proteolysis \citep{Tsuboyama2023MegaScale}. Yet those
data do not remove the need for nanobody-specific target supervision. The
mega-scale domains are 40 to 72 residues long, whereas VHH domains are longer, and
models fine-tuned on such data transfer less well to 177 to 501-residue proteins
than to small domains \citep{Chu2024EsmTherm}. For nanobody Tm prediction,
target-specific data and target-centered validation remain essential
\citep{ValdesTresanco2023NbThermo,Zhang2025NbBench,Schmirler2024FineTuning}.

Here, we test whether computed stability labels can supply useful supervision
when experimental nanobody Tm labels are scarce. We treat Tm prediction as the
target task and each computed label as a computational task that shares an ESM2-based
sequence representation \citep{Lin2023Evolutionary}. The computational labels include
mutation free energies from alchemical free-energy perturbation (FEP), Rosetta
mutation scores, ThermoMPNN stability scores, Rosetta scores on variants
proposed by ESM2 or sampled at random, and an MD-derived native-contact stability
label computed either over the same mutation scans used for FEP or, as a control,
over a diverse panel of nanobodies. To keep the comparison target-centered and
fair across very different labels, every source is tuned independently and selected
on experimental Tm validation data, and every final claim is evaluated on the same
held-out experimental Tm test set with paired error statistics.

\begin{figure*}[t]
\centering
\includegraphics[width=0.66\textwidth]{figure_1_concept_protocol_v2}
\caption{\textbf{Multi-task transfer-learning architecture for using simulation
data in experimental Tm prediction.}
Experimental sequence data (left) and computational sequence data (right) are
encoded by a shared ESM2 protein language model into per-sequence embeddings.
The embeddings pass through a shared multilayer perceptron and then split into
two task-specific heads: a Tm head that predicts the experimental melting
temperature, and a computational head that predicts the computational scores
(mutation $\Delta\Delta G$ values or an MD-derived Q-value). Only the
experimental Tm head defines the final objective, while the computational head lets the
simulation-derived labels reshape the shared representation during training.
Flame and snowflake icons mark trainable and frozen parameters, respectively;
the ESM2 encoder is either fine-tuned (hot) or kept frozen depending on the
configuration.}
\label{fig:learning_framework}
\end{figure*}

The results are organized by two design axes. First, \emph{data design}: the same
MD native-contact observable does not improve Tm when computed over a
diverse-nanobody screen, but does when computed on the FEP-matched mutation scans,
where it matches FEP with a frozen encoder --- so matching the auxiliary source's
sequence design to the target, not the choice of dynamical quantity, is what lets
an MD label transfer.
Second, the \emph{physical observable}: the alchemical FEP free energy improves Tm
whether the encoder is frozen or fine-tuned, whereas the native-contact label helps
only a frozen encoder, so the physical content of the label sets how deeply it
transfers. The other structure-based scores are weak or null under the same
tuning. These results extend the Sim2Real argument to a protein-engineering setting
in which the computational and experimental labels differ in physical meaning, and
turn "run more simulation" into two concrete design choices --- matched data design
and a transferable observable --- that decide whether computed data help.
